% ID: SEL-D0078-OA-E01
% ESTADO: cotejada. Segunda pasada visual y borrador compilado revisados.
% Fuente: Selectividad/sel_2008_matematicas_aplicadas/Mat CCSS 6 - Examen y criterios de correción.pdf
% SHA-256: 256ad5f6dc324a5bd081279a0455d3ea64e9feb86c8be52341faa83024bc3b5b
% Página: 1; opción A; ejercicio original 1.

\textbf{(3 puntos)} Una empresa produce botellas de leche entera y de leche desnatada y tiene una capacidad de producción máxima de 6000 botellas al día. Las condiciones de la empresa obligan a que la producción de botellas de leche desnatada sea, al menos, la quinta parte de las de leche entera y, como máximo, el triple de la misma. El beneficio de la empresa por botella de leche entera es de 20 céntimos y por botella de leche desnatada es de 32 céntimos. Suponiendo que se vende toda la producción, determine la cantidad de botellas de cada tipo que proporciona un beneficio máximo y el importe de este beneficio.
